\section{Introduction}\label{sec:intro}
The Weil criterion expresses the Riemann hypothesis as positivity of an
explicit quadratic form on compactly supported test functions. An inverse
bulk approach asks for a quantum system, a boundary or defect preparation,
and an independently positive physical pairing whose pullback is exactly
that form. Interactions are compatible with this objective: a preparation
linear in the input produces a quadratic state norm even when the
underlying dynamics is interacting. What must be derived is the entire
pairing, with its physical adjoint and normalization.

Here ``boundary'' means the locus or observable sector through which
the arithmetic input enters the system. It may be a spatial boundary,
a line defect, an interface, a gluing cut, or a module for a protected
operator algebra. This flexibility permits the use of sphere and Schur
quantization, where explicit representations and positive pairings are
available \cite{Sphere,Schur,RG}. Conformal symmetry can constrain a chosen
defect, but does not identify arbitrary deformations or source maps.
Pure four-dimensional Yang--Mills is also a possible long-term construction
target; no identification with it is made in this paper.

\clearpage
\subsection{Principal result and its limitation}
The most concrete arithmetic result concerns the elementary Schur
representation. For a fixed quantization parameter $0<q<1$, we construct
an electric dressing and one magnetic insertion whose norm on Laurent
polynomials $h$ is
\begin{equation}\label{eq:intro-norm}
 \norm{\Phi_{q;r,d}[h]}^2
 =\int_{S^1}|h(z)|^2
 \frac{dr(1+r)}{1-r}\frac{|1-z|^2}{|1-rz|^2}\frac{\dd\theta}{2\pi},
 \qquad z=e^{i\theta},\quad 0<r<1,
\end{equation}
with $d>0$. Its Fourier coefficients are
\begin{equation}\label{eq:intro-coeff}
 \frac{dr(1+r)}{1-r}\frac{|1-z|^2}{|1-rz|^2}
 =\frac{2dr}{1-r}-d\sum_{n\geq1}r^n(z^n+z^{-n}).
\end{equation}
Thus $r=p^{-1/2}$ and $d=\log p$ give every required single-prime
repetition coefficient. The arithmetic parameters enter the preparation;
$q$ is held fixed. Theorem~\ref{thm:schur-prime} proves this identity and
the domain statement needed for its magnetic insertion. An explicit
unitary decomposition then relates the circle calculation to the real
logarithmic coordinate of the Weil form.

The positive contact $2dr/(1-r)$ is part of the result. This is a Schur
realization of the positive prime reference already present in the
project's conservative return construction. Neither the reference form
nor the geometric Fourier identity is claimed as new. The additional
content is the specified operator preparation, its source normalization,
its graph-domain justification, and its compatibility with one fixed
Schur parameter. A physical interface implementing the continuous-input
lift has not been constructed.

\subsection{Other results and the remaining question}
The Gaussian real boundary module identifies the archimedean
quarter-shift and oscillator levels, but its natural state-smearing norm
has a different spectral weight. In the interacting $T[SU(2)]$ sphere
model we evaluate an infinite family of charged-state norms and show
that bare powers do not have geometric return ratios. A natural electric
preparation in the conformal $SU(2)$ theory with four hypermultiplets
also fails, as does the undressed elementary Schur return prescription.
These are exclusions of specified preparations, not of the theories.

Two further results govern the new dressed source. Its finite magnetic
domain does not permit unrestricted use of higher-power identities.
Moreover, gluing the prime filters through fixed output vectors produces
an atom at $\log(3/2)$ unless those vectors are orthogonal
(Proposition~\ref{prop:mixed}). Orthogonality restores the independent
positive contacts. A successful joint construction must therefore change
the preparation or the gluing operation in a substantive way.

The sharpest result of this version says how far that change must go.
Section~\ref{subsec:interference} records the elementary bound
$\norm{(G+S)f}^2\geq(\norm{Gf}-\norm{Sf})^2$ and applies it to the two
positive channels that the dressed construction produces, the gamma energy
and the summed prime references. The bound fails at explicit test functions:
at $L=5/4$ the indicator of the interval already violates it, and it fails
by a factor tending to infinity as $L$ grows. Hence no coherent source has
those two forms as its channel norms (Theorem~\ref{thm:two-channel}), even
though each is realized exactly and every prime coefficient matches. The
contact and the pole term have to come from a compression or a projection
inside one space, not from an additional positive channel
(Proposition~\ref{prop:no-augment}). Section~\ref{subsec:spectral} records
the spectral form of the target, which explains the mechanism: the quantity
to be matched is smaller than each term being matched by several orders of
magnitude.

Section~\ref{subsec:subtraction} then supplies such a presentation. Beside
the positive prime reference there is a mirror reference
\eqref{eq:mirror-weight} carrying the same prime atoms with the opposite sign
and a strictly smaller contact $2dr/(1+r)$. Its atoms cancel the Weil prime
term outright, so the target becomes a difference
(Theorem~\ref{thm:subtraction}) in which the remaining form $A_L$ has no
prime translation at all, and the Weil criterion on $I_L$ becomes one
domination statement, equivalently a compression
(Corollary~\ref{cor:domination}). The prime-free form is not marginal where
$Q_L$ is, which is what makes it a construction target rather than a
restatement.

Two results of this version simplify that target and withdraw a claim of the
previous one. Section~\ref{subsec:reachable} characterises the states the
elementary magnetic insertion can reach: exactly the functions analytic on a
disk of radius exceeding one. The factor the operator inserts therefore
constrains nothing, the mirror reference is realized by one dressing at the
same fixed $q$ (Corollary~\ref{cor:mirror-schur}), and producing a prescribed
weight in a positive sector is not by itself evidence of arithmetic structure
in the representation. Section~\ref{subsec:gammacompression} then splits the
pole form by parity and presents the target as $Q_L=K_+-T_L$, with $K_+$ the
gamma energy together with a constant and one rank-one term --- a norm, with an
explicit source --- and $T_L$ the negative contact, the odd half of the pole
form and every mirror reference, positive term by term. Weil positivity on
$I_L$ becomes the single domination $T_L\preceq K_+$. Neither side needs an
unproved inequality, so the prime-free inequality is not a prerequisite for the
construction problem; and by Corollary~\ref{cor:spectral} the pole term never
required an indefinite-metric sector, which the previous version asserted in
Section~\ref{subsec:leaves} and which is withdrawn here.

The background and normalization are developed in
Section~\ref{sec:weil}, independently of any companion manuscript.
Section~\ref{sec:positive} states the conditional construction-to-RH
implication and a general source-regularity restriction.
Sections~\ref{sec:sphere}--\ref{sec:gluing} contain the main calculations.
Section~\ref{sec:outlook} formulates the next complete fitting problem.
The appendices retain positive gamma refinement and earlier gauge and
rational-feedback tests, with proofs sufficient to use them without
consulting the research notes.

The classical Weil criterion, gamma identities, and established
representation formulas are explicitly attributed. The comparison
calculations are working results, not a claim of literature priority.
In particular, an exact identity between local phase derivatives and
the prime terms is not a positivity theorem. No candidate in this draft
matches the complete first-prime Weil form, and no all-support result
follows from the local norm in \eqref{eq:intro-norm}.
